\input{../tex-inputs/latex-leseansicht-vorspann}

\section[Arthur Schnitzler an Gustav Schwarzkopf, {[}13. 8. 1924{]}]{L04524 Arthur Schnitzler an Gustav Schwarzkopf, {[}13. 8. 1924{]}}
\nopagebreak\mylabel{L04524v}
\rehead{ }\normalsize\beginnumbering\briefempfaengerindex{Schwarzkopf, Gustav@\textsc{Schwarzkopf, Gustav}!zzzSchnitzler, Arthur@\emph{von Arthur Schnitzler}!1924-08-131@{{[}13. 8. 1924{]}}|(be}
\toendnotes[C]{\smallbreak\pagebreak[2]}
\correspDesc{Versand  durch Arthur Schnitzler am [13. 8. 1924] in Celerina
\newline{}Übermittlung  am 13. 8. 1924 in Celerina
\newline{}Erhalt  durch Gustav Schwarzkopf im Zeitraum [14. 8. 1924
                  – 18. 8. 1924?] in Wien}\toendnotes[C]{\smallbreak}
\Standort{DLA, A:Schnitzler, HS.1985.1.1897.}
\physDesc{Bildpostkarte, 638 Zeichen
\newline{}Handschrift: Bleistift, lateinische Kurrent
\newline{}Versand: Stempel: »\nobreak{}\oindex{Celerina@\textbf{Celerina}|pwk}Celerina Graubünden, 13. VIII. 24, 11\nobreak{}«.  }\toendnotes[C]{\smallbreak}
\pstart
           \noindent{}\centering{}{\pb}\textcolor{gray}{\textbf{Celerina}}\oindex{Celerina@\textbf{Celerina}|pw}\pend
           \pstart{}{\pb}Hrn.\pend{}\pstart{}Gustav Schwarzkopf\pend{}\pstart{}Wien I\oindex{I., Innere Stadt@\textbf{I., Innere Stadt}|pw}\pend{}\pstart{}Tiefer Graben 17\oindex{Wien@\textbf{Wien}!I., Innere Stadt@\textbf{I., Innere Stadt}!Tiefer Graben 17@\textbf{Tiefer Graben 17}|pw}\pend{}{\bigskip}\vspace{1em}
\pstart
           \noindent{}lieber Guſtav, nun hab ich Wien\oindex{Wien@\textbf{Wien}|pw}
               doch verlassen (am 5. August) ohne Ihnen Adieu zu sagen. So send ich
               Ihnen de{\geminationn} von hier aus viele herzliche Grüße. Ich habe
               mich wieder in dem herrlichen Celerina, Cresta
                  Palace\oindex{Cresta Palace@\textbf{Cresta Palace}|pw}, angesiedelt. Meine Leute\pwindex{Cappellini, Lili 13.\,9.\,1909 Wien – 26.\,7.\,1928 Venedig@\textsc{Cappellini, Lili} (13.\,9.\,1909 Wien – 26.\,7.\,1928 Venedig)|pwv}\pwindex{Schnitzler, Olga 17.\,1.\,1882 Wien – 13.\,1.\,1970 Lugano@\textsc{Schnitzler, Olga} (17.\,1.\,1882 Wien – 13.\,1.\,1970 Lugano)|pwv}\pwindex{Schnitzler, Heinrich 9.\,8.\,1902 Hinterbrühl – 12.\,7.\,1982 Wien@\textsc{Schnitzler, Heinrich} (9.\,8.\,1902 Hinterbrühl – 12.\,7.\,1982 Wien)|pwv}  in Zuoz\oindex{Zuoz@\textbf{Zuoz}|pw} (eine halbe Stunde von hier.) Heini\pwindex{Schnitzler, Heinrich 9.\,8.\,1902 Hinterbrühl – 12.\,7.\,1982 Wien@\textsc{Schnitzler, Heinrich} (9.\,8.\,1902 Hinterbrühl – 12.\,7.\,1982 Wien)|pw} fährt am Samstag nach Wien\oindex{Wien@\textbf{Wien}|pw}, am 25. 8. nach Berlin\oindex{Berlin@\textbf{Berlin}|pw}. Es ist möglich, dß ich in 10  Tagen etwa von hier an
               den Vierwaldstättersee\oindex{Vierwaldstättersee@\textbf{Vierwaldstättersee}|pw}, ins Berner Oberland\oindex{Berner Oberland@\textbf{Berner Oberland}|pw} oder an einen oberitalienischen See\oindex{Oberitalienische Seen@\textbf{Oberitalienische Seen}|pw} gehe. Woraus Sie {\pb}ersehen können, daſs \label{K_L04524-1v}\edtext{Amerika\oindex{Vereinigte Staaten von Amerika [USA]@\textbf{Vereinigte Staaten von Amerika [USA]}|pw}\orgindex{Thomas Seltzer, Inc.@Thomas Seltzer, Inc.|pw}}{\lemma{\textnormal{\emph{Amerika}}}\Cendnote{\textnormal{Im \emph{Tagebuch}\pwindex{Schnitzler, Arthur 15. 5. 1862 Wien – 21. 10. 1931 ebd.@\textsc{Schnitzler, Arthur} (15. 5. 1862 Wien – 21. 10. 1931 ebd.)!Tagebuch@\strich\emph{Tagebuch}|pwk} klagte Schnitzler:
                     »Meine finanz. Verhältnisse verschlechtern sich rapid; dadurch dass die
                     Zahlungen von allen Seiten ausbleiben. Insbesondre Amerika\oindex{Vereinigte Staaten von Amerika [USA]@\textbf{Vereinigte Staaten von Amerika [USA]}|pw}, hier vor allem der Verleger Thomas Seltzer\pwindex{Seltzer, Thomas 22.\,2.\,1875 Poltava – 11.\,9.\,1943 New York City@\textsc{Seltzer, Thomas} (22.\,2.\,1875 Poltava – 11.\,9.\,1943 New York City)|pw} beträgt sich als Gauner niedrigster
                     Sorte.«, siehe A. S.: \emph{Tagebuch}, 21. 7. 1924.}}}\label{K_L04524-1} noch immer nicht zahlt{[}.{]}\pend
           
\pstart
           Und von \label{K_L04524-2v}\edtext{Herterich\pwindex{Herterich, Franz 3.\,10.\,1877 München – 28.\,10.\,1966 Wien@\textsc{Herterich, Franz} (3.\,10.\,1877 München – 28.\,10.\,1966 Wien)|pw}}{\lemma{\textnormal{\emph{Herterich}}}\Cendnote{\textnormal{Franz Herterich\pwindex{Herterich, Franz 3.\,10.\,1877 München – 28.\,10.\,1966 Wien@\textsc{Herterich, Franz} (3.\,10.\,1877 München – 28.\,10.\,1966 Wien)|pwk} war seit 1923
                  Direktor des \emph{Burgtheaters}\orgindex{Burgtheater@Burgtheater|pwk}, wo für den
                     Herbst die Uraufführung von \emph{Komödie der Verführung}\pwindex{Schnitzler, Arthur 15. 5. 1862 Wien – 21. 10. 1931 ebd.@\textsc{Schnitzler, Arthur} (15. 5. 1862 Wien – 21. 10. 1931 ebd.)!Komödie der Verführung. In drei Akten@\strich\emph{Komödie der Verführung. In drei Akten}|pwk}\eventindex{Burgtheater@\textbf{Burgtheater}!Uraufführung von Komödie der Verführung, 11.10.1924@Uraufführung von Komödie der Verführung, 11.10.1924|pwk} geplant war und in derselben Spielzeit auch die Wiener\oindex{Wien@\textbf{Wien}|pwk}{ }Erstaufführung\eventindex{Burgtheater@\textbf{Burgtheater}!Premiere von Der Schleier der Beatrice, 23.5.1925@Premiere von Der Schleier der Beatrice, 23.5.1925|pwkv} von \emph{Der Schleier der Beatrice}\pwindex{Schnitzler, Arthur 15. 5. 1862 Wien – 21. 10. 1931 ebd.@\textsc{Schnitzler, Arthur} (15. 5. 1862 Wien – 21. 10. 1931 ebd.)!Schleier der Beatrice. Schauspiel in fünf Akten@\strich\emph{Der Schleier der Beatrice. Schauspiel in fünf Akten}|pwk} stattfinden
                  sollte.}}}\label{K_L04524-2} hör ich nichts. – Beste Grüße an Doktor Max\pwindex{Schwarzkopf, Max 12.\,6.\,1857 Wien – 14.\,4.\,1928 ebd.@\textsc{Schwarzkopf, Max} (12.\,6.\,1857 Wien – 14.\,4.\,1928 ebd.)|pw}.\pend
           \pstart Herzlichſt wie immer Ihr \spacefill\mbox{A.}\pend{}\selectlanguage{ngerman}\endnumbering\briefempfaengerindex{Schwarzkopf, Gustav@\textsc{Schwarzkopf, Gustav}!zzzSchnitzler, Arthur@\emph{von Arthur Schnitzler}!1924-08-131@{{[}13. 8. 1924{]}}|)be}\mylabel{L04524h}
\begin{anhang}
\end{anhang}\newcommand{\dateiname}{L04524}\newcommand{\titel}{Arthur Schnitzler an Gustav Schwarzkopf, [13. 8. 1924]}\newcommand{\editorInnen}{Herausgegeben von Jahnke, SelmaMüller, Martin Anton}\input{../tex-inputs/latex-leseansicht-abspann}
